% META: {"id": "TSPE-LOG-EX-012", "chapitre": "TSPE-LOGARITHME", "type_objet": "exercice", "capacites": ["TSPE-LOGARITHME-C2"], "capacites_codes": ["C2"], "methodes": ["M1"], "parcours": 2, "competences": ["modeliser", "chercher", "calculer"], "duree_min": 8, "mode_creation": "ex_nihilo", "sources_inspiration": [], "fichier_tex": "chapitres/TSPE-LOGARITHME/exercices/TSPE-LOG-EX-012.tex", "corrige_tex": "chapitres/TSPE-LOGARITHME/corriges/TSPE-LOG-CO-012.tex", "status": "approved"}
% BEGIN-VERIFY
% from sympy import Rational, ln, log
% n_exact = ln(2) / ln(Rational(103, 100))
% assert round(float(n_exact), 2) == 23.45
% import math
% assert math.ceil(float(n_exact)) == 24
% END-VERIFY
\begin{exercice}{TSPE-LOG-EX-004}{3}{18}
Un capital de $1000$ euros est place a interets composes au taux annuel de $3\%$. Le capital au bout de $n$ annees est $C(n) = 1000 \times 1{,}03^n$.
\begin{enumerate}
  \item Determiner, en resolvant une equation avec $\ln$, la valeur exacte du reel $n$ pour lequel $C(n) = 2000$ (le capital double).
  \item En deduire le plus petit entier $n$ a partir duquel le capital a double (justifier en utilisant la croissance de la suite $(C(n))$).
\end{enumerate}
\end{exercice}
